\subsection{Guarded finite-state controller}\label{sec:method-m01}
The guarded finite-state controller (M01) is an explicit transition policy over the observable task state and the currently registered, admissible candidates. It does not fit weights, estimate confidence, or generate action arguments. The policy reads the state flags published by the environment, forms the set of operation kinds available among admitted candidates, then applies an ordered rule list. The rule order is part of the method: stale observations are refreshed before duplicate events are acknowledged; missing dependencies are resolved before cancellation or work; cancellation requests precede ordinary workflow continuation; validation precedes acceptance or quarantine; uncertain HTTP writes are reconciled before retry; and conflicts cause rebase before compare-and-swap. For data workflows it advances extract, transform, and report according to explicit completion flags. If no supported rule has an admitted candidate, it stops rather than inventing a new operation.

Let $s_t$ be the visible state record, $A_t$ the registered candidates, and $G(a,s_t)$ the broker's current admission predicate. The operation set and M01 choice are
\begin{equation}
O_t=\{\operatorname{op}(a):a\in A_t,\ G(a,s_t)=1\},\qquad
o_t^*=\operatorname{first}_{\prec}\{o:R_o(s_t)=1\land o\in O_t\},
\end{equation}
where $R_o$ is the hand-authored guard for operation $o$ and $\prec$ is the fixed priority order. Candidate selection is then restricted to matching admitted candidates:
\begin{equation}
a_t=\arg\min_{a\in A_t}\bigl(\mathbf{1}[G(a,s_t)=0]+\mathbf{1}[\operatorname{op}(a)\ne o_t^*]\bigr),
\end{equation}
with deterministic first-candidate tie resolution in the implementation. If there is no $o_t^*$, the result is a typed stop decision with reason \texttt{no\_supported\_rule}; it is not converted to a no-op effect.

The controller's cost is linear in the number of guards and candidates for a direct scan, $O(|R|+|A_t|)$, with constant memory beyond the current observation. There is no fitting split, random seed, calibration set, or checkpoint. Its reproducible parameters are the ordered guard definitions and the operation mapping in the pinned `reflexmesh` package; changing guard priority is a method change and requires a new source digest and evaluation. This method can be preferable when process obligations are stable and a state flag has an exact operational meaning. It is brittle when the state vocabulary is incomplete, two rules conflict, or a new task needs semantic interpretation. Rule coverage should be assessed against independently authored states, including unknown flag combinations and empty candidate sets, rather than only the happy path used to author each rule.

M01 is a strong baseline on this benchmark because many task families expose exactly the flags and transitions the rules encode. It solves 1,200 of 1,200 core episodes at 0.0219 ms median and 0.0394 ms p95 decision latency in the frozen run. This result is evidence about this finite operation suite, not about general software automation. It does not prove completeness under production concurrency, establish external validity, or make the rule set a universal optimum. An independent postcondition remains necessary: a correct rule choice can still encounter a failed effect, changed resource, or unresolved external write. The acceptance test therefore checks both that every selected candidate passed broker admission and that the final goal predicate is true; admission alone is not success.



\paragraph{Additional implementation analysis.} The implementation is intentionally inspectable at the guard level. A reviewer can enumerate the rule order and identify which observable field activates each branch; no latent state or fitted checkpoint obscures this mapping. This advantage depends on observation quality. If an event is missing, delayed, or represented by a value outside the guard vocabulary, the policy can take a reasonable-looking but wrong branch or stop. Consequently, state freshness and the existence of a candidate are preconditions, not incidental metadata. A test of this controller should enumerate each rule's true and false boundary, each pair of simultaneously true guards, empty candidate lists, stale revisions, and repeated execution after an effect record. The benchmark's 100% aggregate completion only establishes that the implemented rules cover the authored transitions; it cannot stand in for those boundary tests.

The action expression above uses a deterministic tie rule over the candidate list as presented by the runtime. Candidate order is therefore part of reproducibility and must be recorded or canonicalized. The admission predicate is evaluated again at the effect boundary, so a candidate can pass policy selection and later be rejected after a revision changes. The controller should then observe the rejection and recompute; retrying the same serialized decision would violate the state-dependent guard. The implementation has no explicit estimate of expected value, risk, or planner cost. That simplicity makes its cost low and its failure modes legible, while placing the burden of coverage on manually maintained rules. A held-out family with renamed actions or partial observations would be a meaningful test of whether the baseline's performance transfers beyond the fixture vocabulary.

\subsection{Behavior-tree controller}\label{sec:method-m02}
M02 represents reactive workflow logic as a behavior tree rather than reusing M01's rule interpreter. The implementation builds a selector root from task-specific sequence nodes. A condition node returns success or failure from the current observation. An action node returns running when a matching admitted candidate exists and failure when it does not. Sequence nodes stop at the first child that is not successful; selector nodes stop at the first child that is not a failure. This status propagation distinguishes an unavailable operation from a completed goal. Each tick emits a node trace containing the node name and status, allowing a reviewer to see which branch was considered and why later siblings were skipped.

For node $n$ and observation $o$, let $\sigma_n(o)$ be its status in $\{S,F,R\}$ for success, failure, and running. An action leaf for operation $q$ has
\begin{equation}
\sigma_{q}(o)=\begin{cases}R,&\exists a\in A(o):\operatorname{op}(a)=q\land G(a,o)=1,\\F,&\text{otherwise.}\end{cases}
\end{equation}
The sequence and selector combinators use short-circuit status rules; a sequence propagates its first $F$ or $R$, while a selector propagates its first $S$ or $R$. For a sequence with children $n_1,\ldots,n_k$, success requires every child to succeed:
\begin{equation}
\sigma_{\mathrm{seq}}(o)=S\iff\bigwedge_{i=1}^{k}\sigma_{n_i}(o)=S,\qquad
\sigma_{\mathrm{sel}}(o)=F\iff\bigwedge_{i=1}^{k}\sigma_{n_i}(o)=F.
\end{equation}
The selected action is the first running leaf reached by the traversal; the broker independently checks it before execution.

The tree has explicit priority branches for stale observation refresh, duplicate acknowledgement, dependency waiting, and cancellation, followed by task subtrees for validation, testing, HTTP, restart, cleanup, extraction, compare-and-swap, jobs, reconciliation, and untrusted content. Its traversal cost is $O(N)$ in the visited nodes and it stores no learned or persistent hidden state. Unlike M01, the tree's hierarchy lets operation-specific conditions be composed and recorded as named nodes. It still uses hand-authored predicates and fixed ordering; hierarchy does not itself learn a better policy or prevent a poorly ordered selector from starving a lower-priority branch. A running leaf is a proposal for an admitted action, not proof that the environment effect finished.

The distinction between tree status and task outcome is material. `running` means that the node found a candidate requiring execution; after execution, a new observation reticks the tree. A leaf that returns failure because no matching candidate is present can cause a selector to consider another subtree, so candidate absence is not silently treated as task failure. Conversely, a sequence can remain running while an action is pending; repeated ticks must not repeat an un-idempotent effect. The runtime's intent identity and effect record provide that boundary, while the tree records the decision path. Cancellation and asynchronous completion require explicit state updates; the behavior-tree abstraction does not imply that cancellation terminates an external process.

M02 is evaluated through the same native broker and independent task postconditions as other policies. Its performance on the core suite is perfect in the archived aggregate, but the suite's deterministic state flags and known action vocabulary favor explicit reactive control. Its proper use is a compositional deterministic baseline for workflows where preconditions and recovery branches are known. Evaluation should expose rule conflicts, repeated events, unavailable leaves, livelock, and starvation, and should count effect retries separately from tree ticks. The benchmark does not establish that a behavior tree is universally more interpretable or robust than a finite-state implementation; it establishes only that this independently implemented traversal meets the current fixture contracts.



\paragraph{Additional implementation analysis.} The tree implementation separates control-flow composition from task-specific predicates. Each leaf identifies an operation kind; the candidate resolver selects a currently admitted candidate for that leaf. The tick trace is useful evidence about evaluated branches, but its presence does not establish that every semantically relevant branch was considered. A compact tree can still encode the same assumptions as a long rule list, and selector ordering can starve lower branches if an earlier branch remains perpetually running. Tests should therefore cover nonterminal results, repeated ticks, branch starvation, missing leaves, and effect record-driven transitions. A tree status also has a narrower meaning than a workflow outcome: a running action node says that work is available, not that it completed.

The method is reset between episodes and has no persistent learned parameters. This makes seed changes irrelevant to its decision function when the same observation and candidate order are supplied. The tree's effective parameters are its serialized topology, conditions, operation names, and ordering; these should be content-addressed as part of lineage. If task authors modify a branch after examining test outcomes, the test family is no longer independent. M02's strong core success therefore provides a deterministic comparator, not evidence that behavior trees are intrinsically robust. A useful stress suite would perturb event order, remove a candidate after a pending status, and inject a conflicting urgent signal while preserving the same task postcondition. Runtime authority still belongs to the broker, not the tree traversal.

\subsection{Logistic imitation policy}\label{sec:method-m03}
M03 is a fitted candidate scorer trained to rank the demonstrated action among candidates that the environment has already registered. Each candidate is represented by the fixed feature vector declared by `reflexmesh.features`; observation-level features and candidate descriptors are combined in a common feature matrix. A binary logistic model estimates demonstration preference, and an independent Platt-style logistic calibrator maps its raw margin to a probability-like score on a disjoint calibration split. The declared target is the verified demonstration action, not task completion, authorization, or future reward. The policy selects the highest-scoring currently admissible candidate and returns its action identifier, scores, target description, and a diagnostic that explicitly says the score is not a safety bound.

For feature vector $x$ and demonstration label $y\in\{0,1\}$, the fitted linear margin and probability are
\begin{equation}
z(x)=w^\top x+b,\qquad p(y=1\mid x)=\operatorname{sigmoid}(z)=\frac{1}{1+e^{-z}}.
\end{equation}
The training run used 12,576 candidate rows; logistic regression used $C=5$, the `liblinear` solver, a 600-iteration cap, CPU execution, and random seed 17. Calibration used 2,136 separate candidates, logistic calibration with $C=10$ and a 300-iteration cap. At inference the choice is
\begin{equation}
a_t=\arg\max_{a\in A_t^{\mathrm{adm}}}\widehat p_{\mathrm{cal}}(y=1\mid x(o_t,a)),\qquad
\text{stop if }A_t^{\mathrm{adm}}=\varnothing.
\end{equation}
This is a ranking decision under a fixed set of choices; the maximum calibrated demonstration score is not a calibrated probability that the resulting task will succeed.

The CPU fitting is inexpensive compared with recurrent or planner inference, and scoring $m$ candidates with $d$ features costs $O(md)$. The exported linear artifact stores ordered feature names, coefficients, intercept, and calibration coefficients. The native implementation loads that explicit schema and rejects a feature-order mismatch rather than silently applying weights to different meanings. Python and Rust parity is checked on frozen vectors, including score comparison. The model contains no episode memory: reset is a no-op, so delayed events are represented only if their information appears in the current observation features. Distribution shift in candidate descriptions, missing feature values, or correlations between demonstration style and task identity can therefore alter ranking without a warning.

The calibration report for this fit uses 2,136 rows and records a Brier score of 0.000040 and log loss of 0.000582 for the demonstration-preference target. These values are not task-success metrics and may be optimistic under the particular split construction; they do not establish that the model is calibrated on transfer families. M03 reaches 1,200/1,200 core successes in the frozen evaluation, but that outcome is compatible with the suite being solvable by fixed rules and does not show that learning contributes beyond M01. The proper ablation is a matched fixed policy with the same candidate set and postcondition oracle. Use M03 when modest-data tabular ranking is useful and the feature contract is stable; use abstention or a reviewed fallback when candidates lie outside the training support.



\paragraph{Additional implementation analysis.} The training corpus contains 12,576 candidate rows, with a separate calibration sample of 2,136 rows. The configured logistic fit uses inverse regularization $C=5$, the liblinear solver, seed 17, and a 600-iteration cap; calibration uses $C=10$ and at most 300 iterations. These settings define a candidate-ranking model for the recorded feature representation, rather than a direct policy over arbitrary natural-language instructions. Row-level randomization would leak near-duplicate candidates from the same task into both fit and evaluation, so reported generalization must be interpreted in light of the experiment's group split and authored-family overlap.

For a multicandidate decision, the scorer converts candidate features into class scores, masks candidates not admitted by the runtime, and selects the highest scoring remaining candidate. The calibrated probabilities reported for a demonstration label are not probabilities that the chosen action will succeed. The measured Brier score of 0.000040 and log loss of 0.000582 concern that label target only. A low score can coexist with poor action utility when the label reflects the authored demonstrator or when labels are nearly deterministic. Calibration should be measured on genuinely held-out task families and against outcome-based utility, with reliability diagrams and subgroup counts. Because M03 has no recurrent history, it cannot use the order of prior observations except where a feature explicitly summarizes it. This method is a compact supervised comparator; it does not estimate the causal value of requesting a planner.

\subsection{Gradient-boosted action policy}\label{sec:method-m04}
M04 is a gradient-boosted decision-tree scorer over the same candidate-feature contract as M03. It fits a binary logistic objective to the demonstrated candidate and uses the tree ensemble's raw logit as an action-ranking score. The training pipeline enables categorical features and uses histogram-based split finding. A separate calibration split fits a logistic mapping over the raw margin. At inference, the policy scores only the candidates supplied in the observation and already admitted by the runtime; it does not generate new tool arguments or override the broker. This makes the comparison with M03 interpretable as a change in scoring family under a largely shared observation and action interface, though the fitting procedures and inductive biases remain different.

For additive trees $f_j$, the ensemble margin after $T$ boosting rounds is
\begin{equation}
F_T(x)=F_0(x)+\sum_{j=1}^{T}\eta f_j(x),\qquad
\widehat p(y=1\mid x)=\operatorname{sigmoid}(F_T(x)).
\end{equation}
The archived fit uses the XGBoost `binary:logistic` objective, 160 estimators, maximum depth 5, learning rate $0.08$, histogram tree method, categorical feature support, two CPU workers, and seed 17. Calibration uses 2,136 held-out candidate rows and the same disjoint Platt-style calibration procedure as the linear scorer. The policy's choice is
\begin{equation}
a_t=\arg\max_{a\in A_t^{\mathrm{adm}}}\operatorname{Calibrate}(F_T(x(o_t,a))).
\end{equation}
The score is calibrated only with respect to the demonstration label distribution represented in the calibration split; it is not a task-success probability and not an execution permission.

Tree scoring costs approximately the sum of tree depths for each candidate, $O(mTd)$ as a conservative dense bound for $m$ candidates, $T$ trees, and feature dimension $d$. Training is more expensive than the logistic baseline but remains offline. The package exports both the original model and an explicit tree representation with base margin, split feature index, threshold, child identifiers, missing-value route, leaf values, and calibration coefficients. Native scoring follows the serialized missing-value branch; a wrong branch or feature ordering changes decisions even if the model artifact hash is correct. Validation therefore compares raw margins and calibrated probabilities on the same vectors before evaluating task outcomes.

In the frozen core, M04 solves every scheduled task and, on the structural-transfer sample, solves 24/24 versus M12's 21/24. Its success on the current data does not establish broad superiority: the core has finite authored families and the transfer sample is small. The calibration report's low aggregate Brier and log-loss values refer to teacher-action preference. The sparse intermediate-confidence bins in the report also show why overall averages can obscure calibration uncertainty. Tree ensembles can exploit categorical interaction patterns that a linear model cannot represent, but they can memorize task-family identifiers or fail on unseen tool semantics. Feature leakage, family overlap, and the held-out split must therefore be audited independently of native parity. M04 is a strong tabular comparator, not a universal task planner.



\paragraph{Additional implementation analysis.} The XGBoost configuration uses 160 trees of maximum depth five, learning rate 0.08, histogram construction, categorical support, two jobs, and seed 17. The model ranks action candidates using the same candidate interface and relies on the broker mask for admissibility. The fitted tree ensemble can capture nonlinear feature interactions without requiring a recurrent state, but its predictions remain tied to the training distribution and feature schema. The run lineage must include the exact library version, category encoding, feature order, split identity, and checkpoint digest; changing categorical handling can change decisions even if the high-level method name stays constant.

On the small structural-transfer comparison M04 succeeds on all 24 cases while the proposed controller succeeds on 21. This is descriptive evidence for that composition, not proof that boosted trees dominate planners generally. It is especially important to keep structural transfer distinct from semantic transfer: M04 records only two successes among 128 cases in the latter lane. The difference indicates that the two lanes measure different generalization demands. A meaningful follow-up should separate task templates and tool semantics before feature fitting, report performance by family, and vary environment behavior within family. Compute should include model loading, feature construction, and any calibration step. XGBoost supplies a strong learned fast-policy baseline; its performance does not test whether slow reasoning is worth its cost.

\subsection{Recurrent imitation policy}\label{sec:method-m05}
M05 adds an episode-local gated recurrent unit to candidate imitation. The network consumes a compact state-feature vector at each decision, updates a 32-dimensional hidden state, embeds each candidate through a 48-unit tanh layer, concatenates candidate embedding and recurrent context, and returns one logit per candidate through a further 48-unit tanh scorer. A softmax over the current candidate set converts logits into a normalized preference distribution. Hidden state is indexed by goal ID and cleared on reset; one goal cannot silently inherit another task's context. This model represents short event history while keeping action arguments, resource versions, capabilities, and broker authority in explicit runtime contracts.

The recurrence is
\begin{equation}
h_t=\operatorname{GRU}(x_t,h_{t-1}),\quad e_{t,a}=\tanh(W_c\phi(o_t,a)+b_c),\quad
\ell_{t,a}=w_2^\top\tanh(W_1[e_{t,a};h_t]+b_1)+b_2.
\end{equation}
The candidate preference is
\begin{equation}
p_{t,a}=\frac{\exp(\ell_{t,a}/T)}{\sum_{a'\in A_t^{\mathrm{adm}}}\exp(\ell_{t,a'}/T)},\qquad
a_t=\arg\max_{a\in A_t^{\mathrm{adm}}}p_{t,a}.
\end{equation}
Training uses sequence batches of 32, cross-entropy against the teacher action, 25 epochs, Adam at $0.003$, gradient norm clipping at 1.0, and seed 17. Temperature is selected on separate calibration sequences from 41 values on an exponential grid spanning $e^{-2}$ to $e^2$. ONNX export is checked against CPU PyTorch at candidate counts 1, 3, 6, and 7 with maximum absolute error threshold $10^{-4}$.

At each time step the hidden state is carried from the GRU output; padded candidates are masked from loss. Sequences are shuffled by episode for minibatch updates, while within-episode event order is retained. The training target remains teacher action selection rather than task success. This distinction matters because imitation can reproduce a systematic teacher mistake. Calibration on demonstrated-action labels does not correct label bias or certify the downstream postcondition. A candidate absent from training may still receive a large score because softmax normalizes over the supplied candidates.

Inference cost is linear in sequence hidden size and candidate count for fixed architecture; memory is $O(H)$ per active goal for $H=32$. Reset must be called when a goal ends, and a process restart clears hidden state unless an explicit persistence contract is added. The model does not independently validate whether a delayed event belongs to the current goal. It relies on event identity and observation revisions in its inputs. M05 reaches a perfect aggregate on the frozen core matrix, but no measured advantage over M01 follows from that result. Its role is to test whether temporal state changes candidate preference on tasks with informative history. A valid follow-up would compare cases where otherwise identical current observations have different histories and verify that recurrent updates improve postconditions without state leakage.



\paragraph{Additional implementation analysis.} The recurrent scorer encodes each candidate with a 48-unit tanh projection, updates a 32-unit gated recurrent state from the goal-local history, and scores candidates with a 48-unit network. Training uses cross-entropy for 25 epochs with Adam at learning rate 0.003, batches of 32, and gradient-norm clipping at one. Temperature is selected on a separate calibration set using a 41-point grid. Native export parity is checked against the reference network with a maximum-error tolerance of $10^{-4}$. These choices make the recurrent state an explicit method resource: episode resets must clear it, goal boundaries must select the correct history, and replay must provide the same ordered observations.

A recurrent representation can summarize sequential context that a per-step classifier cannot, but the existence of hidden state does not establish useful memory. The training objective still follows demonstrations, not verified task utility. The report should therefore distinguish sequence fit, ranking behavior, outcome success, and transfer. The parity test checks implementation agreement on a finite vector set and one checkpoint; it does not measure calibration after quantization or robustness to long histories. Runtime timing should include the state update and candidate scoring but exclude unrelated planner latency only when that distinction is explicitly reported. A stronger memory test would ablate history length, perturb event order, and use held-out families where the same local state requires different actions because of earlier events. The frozen core results do not isolate such a benefit.


\paragraph{Decision path and complexity.} The evaluation resets hidden state at episode boundaries and feeds observations in recorded event order. This makes the recurrent controller reproducible from the event stream, but exposes it to truncation and ordering assumptions. A production implementation would need explicit goal-keyed state storage and a policy for partial histories; neither follows from one aggregate score. The method cost is candidate encoding plus recurrent state update. A parity tolerance and task success are separate checks. The recurrent hidden vector is useful only if it changes choices on states whose earlier history matters, so a held-out evaluation should perturb event order and compare decisions with and without history on independently authored families.

\subsection{Contextual linear bandit}\label{sec:method-m06}
M06 implements a shared linear LinUCB action-feature model. It projects candidate features into 64 dimensions with a fixed random sign matrix whose entries are $\pm1/\sqrt{64}$. The design matrix starts at identity, the reward vector at zero, and the exploration coefficient at $\alpha=0.2$. After an observed action reward, inverse covariance and coefficient vector are updated with the Sherman--Morrison identity. At decision time each admissible candidate receives predicted linear reward plus an uncertainty bonus. The policy chooses the largest upper-confidence score with deterministic tie-breaking. It is myopic: the target is one-step reward, not total completion or long-horizon utility.

For projected feature $z_t$, inverse covariance $A_t^{-1}$, reward accumulator $b_t$, and $\theta_t=A_t^{-1}b_t$, the score is
\begin{equation}
U_t(a)=z_{t,a}^{\top}\theta_t+\alpha\sqrt{z_{t,a}^{\top}A_t^{-1}z_{t,a}},\qquad
\theta_t=A_t^{-1}b_t.
\end{equation}
For observed feature $z$ and reward $r$, the update is
\begin{equation}
u=A^{-1}z,\quad A^{-1}\leftarrow A^{-1}-\frac{uu^\top}{1+z^\top u},\quad
b\leftarrow b+rz.
\end{equation}
The archived training performs 2,096 updates from executed one-step rewards, with 64 dimensions, seed 17, and identity regularization equivalent to ridge parameter 1.0.

The logged propensity is 1.0 for the selected action because inference is deterministic. That value does not create overlap for unchosen actions; inverse-propensity estimates cannot reconstruct their outcomes from these logs. The policy does not explore randomly during held-out evaluation and receives no counterfactual labels for unexecuted actions. This limits off-policy evaluation and can make early choices self-reinforcing. Candidate projection reduces compute but can alias distinct categorical tool/action combinations.

Prediction costs $O(mdq+q^2)$ for $m$ candidates, source feature dimension $d$, and projected dimension $q=64$. The confidence bonus assumes the linear reward model and feature support are informative. Terminal task success is not the same target as immediate reward, so a high LinUCB score should not be described as likely completion. This is a controlled bandit baseline for recoverable one-step decisions; multi-step delayed effects require a sequential model or a reward definition that includes future consequences. The present suite does not establish that the bandit improves on fixed rules.



\paragraph{Additional implementation analysis.} The contextual bandit projects the 64-dimensional feature vector through a fixed random-sign projection, starts with an identity inverse covariance estimate, and uses exploration coefficient 0.2. Its fitted policy records 2,096 updates. The fit uses seed 17, while the benchmark policy selects deterministically with propensity one after training. This configuration is a myopic contextual learner: it estimates reward for the current action and does not represent how an action changes later state, how a process completes asynchronously, or how a planner might change future options. The feature projection and reward definition are therefore decisive components, not implementation details.

For a LinUCB-style score, uncertainty is based on the projected feature's leverage under the accumulated design matrix. That term encourages exploration during learning, but a deployed deterministic policy with propensity one no longer randomizes. If the benchmark's task rewards are sparse or only observed after a whole trajectory, credit assignment to intermediate selections is a challenge. The update log and reward mapping should be audited for leakage from postconditions into features, and the selection propensity must be retained if any inverse-propensity estimate is later attempted. A contextual-bandit baseline is appropriate when decisions have local reward and limited action-dependent continuation, not as a drop-in replacement for an asynchronous controller. A useful evaluation would compare online regret under randomized safe exploration and log every candidate propensity; current aggregate task success alone cannot show that the learned uncertainty bonus helps.


\paragraph{Decision path and complexity.} At inference, the candidate score combines estimated reward with a leverage-based uncertainty bonus, after which only broker-admitted actions remain eligible. The random projection reduces representation cost but can discard distinctions important to rare tool preconditions. An auditable run therefore needs the projection seed and matrix, reward definition, action propensity, and update order in its training manifest. The 2,096 updates are transitions, not necessarily 2,096 independent tasks. Since deployed selection is deterministic with propensity one, fitting-time exploration must not be attributed to benchmark-time decisions. State-update cost scales with projected dimension, while scoring scales with candidate count. This baseline tests local reward learning rather than a complete long-horizon strategy.

\subsection{Masked proximal policy optimization}\label{sec:method-m07}
M07 trains a feedforward policy with Maskable PPO in the disposable software environment. At every observation it builds a Boolean mask over the padded action space, marks only currently admissible candidates, and passes that mask to the policy distribution. The network is nonrecurrent and uses separate policy and value multilayer perceptrons with hidden widths 64 and 32. The environment's native admission and task checks remain active during training. Held-out inference orders candidates consistently, applies the same mask, requests a deterministic policy action, and maps its index back to an action ID. An empty valid-action set produces a typed stop instead of selecting padding.

For trajectory sample $(s_t,a_t)$, advantage $\hat A_t$, and probability ratio $r_t(\theta)$, PPO's clipped objective is
\begin{equation}
L^{\mathrm{clip}}(\theta)=\mathbb E_t\left[\min\left(r_t(\theta)\hat A_t,
\operatorname{clip}(r_t(\theta),1-\epsilon,1+\epsilon)\hat A_t\right)\right].
\end{equation}
The sampled distribution is restricted to current valid actions by
\begin{equation}
\pi_{\theta}^{\mathrm{mask}}(a\mid s_t)=\frac{\pi_\theta(a\mid s_t)\mathbf{1}[a\in A_t^{\mathrm{adm}}]}
{\sum_{a'}\pi_\theta(a'\mid s_t)\mathbf{1}[a'\in A_t^{\mathrm{adm}}]}.
\end{equation}
The recorded fit completes 16,384 CPU timesteps, seed 29, with rollout 256, minibatch 64, five epochs, Adam rate $5\times10^{-4}$, $\gamma=0.95$, $\lambda_{\mathrm{GAE}}=0.95$, and clip range 0.2. Entropy coefficient is zero, value coefficient 0.5, and maximum gradient norm 0.5.

PPO can optimize delayed reward in a simulator, unlike a contextual bandit's one-step update, but its evidence depends on simulator fidelity and reward shaping. Exploration occurs during training, not against an arbitrary user workspace; held-out evaluation uses a deterministic trained policy. The policy mask prevents selection of candidates marked inadmissible in the observation, but an action may become stale before execution. The native broker rechecks current authority and revisions, so the mask is a learning constraint, not a replacement for authorization.

Training required approximately 840 seconds in the archived run; this offline cost should not be confused with per-decision latency. M07 solves all core episodes in the frozen aggregate, but the saturated suite cannot isolate whether PPO learning is necessary. Evaluation should include independently varied transitions and reward-hacking checks alongside success, cost, invalid-action attempts, and effect records. The clipped update limits certain policy changes on sampled data; it does not establish regret bounds under distribution shift.



\paragraph{Additional implementation analysis.} The policy uses a masked discrete action space so inadmissible candidates are removed from the categorical distribution. The training run uses 16,384 environment steps, seed 29, CPU execution, rollout length 256, minibatches of 64, five optimization epochs, learning rate 0.0005, discount 0.95, GAE parameter 0.95, and clipping 0.2. Policy and value networks have hidden layers of 64 and 32 units. Training took 840 seconds in the recorded environment. The policy is nonrecurrent; its observation features must expose every historical variable needed to select an action. The mask must be applied identically in training and evaluation to avoid invalid-action shortcuts.

PPO optimizes a clipped policy-gradient surrogate, but that objective is not the product's task-success metric. Reward shaping, episode truncation, and the distribution of terminal outcomes determine which behavior is reinforced. To reproduce the result, the report must identify the exact reward function, normalization, wrapper order, action-to-candidate map, and checkpoint. Training compute is not included in per-decision inference latency; both should be reported when judging a system intended to operate under a resource budget. A held-out evaluation should include run-to-run training variance and family-level uncertainty rather than treating deterministic episodes from one checkpoint as independent samples. PPO is a standard sequential-control comparator here; its inclusion does not establish that reinforcement learning is inherently slower, less reliable, or less suitable than supervised methods.


\paragraph{Decision path and complexity.} At inference, the policy computes logits for candidate positions, masks disallowed actions, and samples or selects from the normalized distribution according to evaluation configuration. Reproducibility depends on candidate ordering and whether stochastic sampling remains enabled. A mask mismatch can change the effective support of the policy. Its forward cost grows with candidate feature encoding and the two-layer network; it stores no recurrent state. Reproduction requires environment wrappers, reward scaling, normalization statistics, action mapping, and checkpoint digest in addition to the training seed and step count. Policy masking is a selection constraint rather than an independent authority boundary, so the external broker remains responsible for final admission.

\subsection{Learned transition and bounded lookahead}\label{sec:method-m08}
M08 trains a compact transition model and uses bounded lookahead to rank current candidates. The network maps each candidate feature vector through 64 and 48 ReLU units to predicted next-state fields, immediate reward, and a termination estimate. Training predicts observed transition fields from executed actions. The terminal output is episode termination, which can represent success or failure; it is not a task-success probability. This semantic distinction must be preserved when reading its candidate utility.

For candidate $a$, the model returns $\widehat x'_a$, $\widehat r_a$, and $\widehat d_a$. The implementation adds one bounded expansion:
\begin{equation}
Q_1(o,a)=\widehat r_a+0.8(1-\operatorname{clip}(\widehat d_a,0,1))
\max_{a'\in A(o'_a)}\widehat r_{a'}.
\end{equation}
Its predictive training objective is a weighted sum of state, reward, and terminal-head errors:
\begin{equation}
L(\phi)=\frac{1}{N}\sum_{i=1}^N\left(\|\widehat x'_i-x'_i\|_2^2+
\lambda_r(\widehat r_i-r_i)^2+\lambda_d(\widehat d_i-d_i)^2\right).
\end{equation}
The run uses Adam at 0.003, 40 epochs, batch size 128, seed 31; horizon is 2 and expansion budget is 24. Validation MSE is 0.00711 overall, with state, reward, and terminal errors persisted separately.

Lookahead regenerates candidate scores under a hypothetical state constructed from predicted fields; it does not enumerate arbitrary plans or model every effect. If the terminal estimate is wrong, the continuation multiplier can overvalue a trajectory. If the transition model is inaccurate, search can exploit its errors. Limiting expansions bounds compute but does not prevent model exploitation. Scores are labeled utility, not calibrated probability, and predicted effects are exposed as diagnostics for the selected action.

Inference work is bounded by the configured candidate expansion loop rather than a full breadth-first tree. Predictive accuracy does not by itself bound policy regret: correlated error can preserve a good ranking, while a small error near two candidate scores can reverse it. The method is useful as a model-based control comparator when transitions repeat and alternatives have delayed consequences. To establish value, a new task set must vary transitions independently and compare closed-loop outcomes, cost, and prediction errors. The current records do not show that learned lookahead surpasses the guarded finite-state policy.


\paragraph{Additional implementation analysis.} The learned transition model predicts next-state features, reward, and termination; it does not predict verified task success as a separate target. Its network has hidden layers of 64 and 48 ReLU units, is trained with Adam at learning rate 0.003 for 40 epochs with batch size 128 and seed 31, and has validation mean-squared error 0.00711. Search depth is limited to two, with a maximum of 24 expansions. These limits bound inference but also restrict the planner's capacity to represent long action dependencies. A low average feature error does not guarantee correct ranking at decision boundaries, especially when errors concentrate on rare cancellation or external-write states.

The model-based planner rolls candidate actions forward through the learned dynamics and evaluates predicted reward. Because termination and success differ, the evaluator must not treat predicted termination as a successful goal. Model error can compound across the two-step horizon; a confidence penalty or ensemble disagreement could expose some uncertainty, but neither would repair distribution shift by itself. The transition validation split and state-feature scaling must be reported alongside the aggregate MSE. A relevant test is decision-aware accuracy: compare the model's chosen first action with a real environment rollout on held-out states, and stratify by event family. The current task suite does not show that this bounded lookahead improves on M01, and reported model error alone cannot support that claim.


\paragraph{Decision path and complexity.} The rollout procedure expands candidate actions through predicted next states, truncates search at the configured depth and expansion cap, and returns the first action from the best predicted path. This bounds computation but can discard a low-ranked branch whose model error would have made it useful. The transition model and evaluator must share feature scaling and terminal semantics. Predicted termination is not verified goal completion, so an independent evaluator checks the real postcondition after execution. Search cost scales with expanded candidates times model forward cost. A decision-level diagnostic should retain selected action, predicted state, realized state, and error by family; mean-squared validation loss alone cannot show that the model ranks first actions correctly.
